\section{INTRODUCTION}\label{ch:intro}

A risk manager who reports a daily value-at-risk (VaR) and expected shortfall (ES) for a portfolio of stocks, Treasuries and gold has to say how those assets move together in a bad week. The correlation between stock and bond returns has no fixed sign. It was about 60\% in the mid-1990s and about $-60\%$ in the early 2000s \citep{baele2010}, and it switched from positive to negative around 2001 \citep{campbell2020}. It turned positive again in 2022: the correlation of daily SPY and IEF returns was $-0.46$ over 2010 to 2019 and $+0.19$ in 2022, and it stayed near zero in 2023 ($+0.04$) and 2024 ($+0.07$) (own calculation from the data of Chapter~\ref{ch:data}, descriptive only).
A forecast that holds the correlation constant would be too optimistic in the tail on exactly the days when the bond holding stops offsetting equity losses. The risk measure matters to practitioners beyond internal reporting: the Basel Committee's market-risk standard sets capital for internal models on expected shortfall at the 97.5\% level, a 2.5\% tail \citep{bcbs2019}, which is the secondary tail level used here. In the hold-out of this thesis, 2020 to 2025, the correlation between the US equity fund SPY and the Treasury fund IEF over the previous 63 trading days was positive on 44.2\% of days (Chapter~\ref{ch:holdout}).

This thesis asks whether letting the SPY--IEF correlation switch with a market regime improves one-day-ahead VaR and ES forecasts of a fixed portfolio, an equal-weight holding of SPY, IEF and the gold fund GLD. The confirmatory comparison was fixed before I used the test sample: a model in which the SPY--IEF correlation depends on the state of a two-state hidden Markov model fitted to SPY returns, against a model with constant correlation and the same GARCH margins. A wider ladder of models, run on development data only, shows where in the chain of models regime information helps; those results are descriptive.

The economic reasoning comes from the stock--bond literature. \citet{campbell2020} tie the sign of the correlation to a slowly changing macroeconomic state, while \citet{baele2010} find that macroeconomic fundamentals explain little and that market conditions such as liquidity and the variance premium matter more. A two-state volatility regime is a market-based signal, so the test speaks to the second reading and not the first (Chapter~\ref{ch:discussion}). The model also learns the stress-day correlation from the five years before each forecast, and before 2022 that correlation was negative or low (Table~\ref{tab:eras}). The question is therefore whether the model improves forecasts, not whether it reproduces the 2022 positive correlation.

The literature leaves the question open (Chapter~\ref{ch:lit}). Regime switching improves volatility-based risk forecasts \citep{ardia2018}, regime-switching correlation has been proposed \citep{pelletier2006} and tested for densities and VaR mainly on stocks \citep{paolella2019}, and univariate and multivariate models have been compared for expected shortfall without regime switching \citep{fortin2023}. I did not find a study of regime-switching correlation that scores VaR and ES jointly and reports what the test could and could not detect.

Forecasts are scored with the FZ0 loss of \citet{patton2019}, which scores VaR and ES jointly, and models are compared with a Diebold--Mariano test with a heteroskedasticity-and-autocorrelation-consistent standard error. The sample is split into development days, on which I explored and which validate nothing, and a hold-out of 1,508 days, 2020 to 2025, used once.

Switching the SPY--IEF correlation lowered the average loss by 0.0013, with a 95\% interval from $-0.0037$ to $+0.0012$, and the one-sided test did not reject ($p=0.158$). A loss difference of that size is roughly the penalty for scaling both the VaR and ES forecasts by a factor between 2 and 3 per cent from the truth (Section~\ref{sec:sesoi}). The smallest difference the test could detect with 80\% probability was 0.0031, and a true improvement equal to the smallest effect of interest, 0.0020, would have been detected with probability about 48\%. The data are consistent with a gain and with none. The result is a measure of how precisely this sample can speak, and I do not read it as evidence that the effect is absent.

The second finding is the more useful one for a practitioner. Both models breached the 5\% VaR on about 6.8\% to 6.9\% of hold-out days, so the comparison is between two forecasts that fell short in the same way. A RiskMetrics EWMA of the portfolio return, scored post hoc on the same days, breached on 5.2\% and had a loss that cannot be told apart from theirs (Section~\ref{sec:holdout-posthoc}). On the development days, regime-specific volatility lowered the loss against a pooled normal in three of four ETF samples, GARCH margins did better still (significant after adjustment in one sample), and the correlation steps were small beside the differences they could detect.

The contribution is an estimate with an interval and a stated detection limit for a question the literature leaves open, not a better model. Chapter~\ref{ch:lit} reviews the literature, Chapters~\ref{ch:data} to~\ref{ch:eval} set out the data, models and evaluation, Chapters~\ref{ch:results} and~\ref{ch:holdout} report the results, and Chapters~\ref{ch:discussion} and~\ref{ch:conclusion} discuss and conclude.
